\section{Moving a path through a neighboring cell}
\label{bridge:section}

We next estimate the cost of replacing one particle's trajectory while
preserving its two outer endpoints.  The replacement passes through a point
chosen uniformly in its home cell or a face-neighbor cell.  Two estimates
serve different purposes: averaging over many separated endpoints makes the
total collision probability against order $D$ obstacles small, whereas a
single replacement has collision probability tending to zero against each
fixed regular obstacle.

We use the regularity conditions of Definition~\ref{path:regular-path}.
Thus $r=\alpha/D$, $s_0=bD^{-1/3}$, $h_0=D^{-1}$, and $\gamma=1/12$.
Each half of $[0,2t]$ has $m_t=\lceil t/r^2\rceil$ mesh intervals of length
$\tau=t/m_t$.  For a regular path $c$, its lifted diameter on interval $k$
is $rA_k(c)$, where
\begin{equation}\label{bridge:oscillations}
 \max_k A_k(c)\le D^\gamma,
 \qquad
 \sum_{k=1}^{2m_t}(1+A_k(c))^3\le 2C_{\rm reg}m_t.
\end{equation}
Its diameter in each outer time window of length $h_0$ is at most $s_0/8$.
All estimates below hold for each fixed $0<t\le1$ and sufficiently large
$D$, uniformly in $\alpha\in[\alpha_0/2,2\alpha_0]$.  In particular, we
may require
\begin{equation}\label{bridge:scale-restrictions}
 \tau\le h_0/2,\qquad h_0<t/4,\qquad rD^\gamma<1,
 \qquad C\tau/t\le r
\end{equation}
for any fixed geometric constant $C$ used below.

Fix a cell $B_v=v+[-1/2,1/2)^3$.  Suppose that a path's outer endpoints
$y,z$ have local lifts satisfying $y\in B_v$ and $|z-v|\le3$.  A
\emph{target cell} is $B_v$ or one of its six face-neighbor cells.  For
$u$ in a target cell $B$, let $Q_{y,z}^u$ be the law obtained by joining two
independent Euclidean Brownian bridges, of generator $\Delta$, from $y$
to $u$ and from $u$ to $z$, each of duration $t$.  We project this path to
the torus.  Let
\[
 Q_{y,z}^{B}=\int_B Q_{y,z}^u\dd u,
 \qquad
 Q_{y,z}^{0}=\text{the Euclidean bridge from $y$ to $z$ of duration $2t$}.
\]
Both laws are projected to the torus.  Unit cells have volume one, so
$Q_{y,z}^{B}$ is a probability measure.
Write $\mathcal C_v$ for the event that the replacement is regular and
its local lift stays in the closed ball of radius $10$ about $v$.
For a deterministic torus path $c$, define its collision event by
\[
 \mathcal H(c)=
 \{\omega: d(\omega(s),c(s))\le r\text{ for some }s\in[0,2t]\}.
\]
The estimates concern $\mathcal C_v\cap\mathcal H(c)$ under the
unconditioned bridge laws.  No independence between the cutoff event and
the collision event is assumed.

\subsection{Separated endpoints and Gaussian densities}

\begin{lemma}[Averaged Gaussian density]\label{bridge:gaussian-packing}
Let $y_1,\ldots,y_m\in\R^3$ have pairwise distances at least $s_0$, where
$m\ge\ell D/2$ and $s_0=bD^{-1/3}$.  Fix $u\in\R^3$ and
$0<s\le t/2$.  The average of the time-$s$ densities of duration-$t$
Brownian bridges from $y_i$ to $u$ is bounded everywhere by
\begin{equation}\label{bridge:packed-density}
 C\left(1+\frac{1}{D s^{3/2}}\right).
\end{equation}
Here $C$ depends on $b,\ell$, but not on $t,s,D$ or $u$.
\end{lemma}

\begin{proof}
The means are $a y_i+(1-a)u$, where $a=1-s/t\ge1/2$, and the covariance
matrix is $2s(1-s/t)$ times the identity.  The means are therefore
$s_0/2$-separated.  For points $q_i$ separated by $s_0/2$, disjoint balls
of radius $s_0/4$ give
\[
 \#\{i:|q_i-x|\le R\}\le C(1+R/s_0)^3.
\]
Split the Gaussian sum into the ball of radius $\sqrt{s}$ about $x$ and
the annuli with radii $j\sqrt{s},(j+1)\sqrt{s}$, $j\ge1$.  The last
bound, multiplied by the Gaussian decay on each annulus, yields
\[
 \sum_i \exp\!\left(-\frac{|x-q_i|^2}{4s}\right)
 \le C\left(1+\frac{s^{3/2}}{s_0^3}\right).
\]
The covariance lies between $s$ and $2s$ times the identity, so the
averaged density is at most
\[
 \frac{C}{m s^{3/2}}+\frac{C}{m s_0^3}
 \le C\left(\frac{1}{D s^{3/2}}+1\right).
\]
\end{proof}

\subsection{Collision with a deterministic obstacle}

\begin{proposition}[Averaged collision bound]\label{bridge:averaged-hit}
Fix $m\ge\ell D/2$ pairs $(y_i,z_i)$ with $y_i\in B_v$ and
$|z_i-v|\le3$.  Assume that each of the two endpoint lists is
$s_0$-separated.  Let $c$ be a deterministic regular torus path whose
outer endpoints are at distance greater than $s_0$ from the corresponding
endpoints of every label under consideration.  Then, for every target
cell $B$,
\begin{equation}\label{bridge:averaged-estimate}
 \frac1m\sum_{i=1}^m
 Q_{y_i,z_i}^{B}(\mathcal C_v\cap\mathcal H(c))
 \le Ctr+C\frac{rD^{3\gamma}}{D\sqrt{h_0}}.
\end{equation}
The same conclusion holds if $c$ is the original path of one of these
labels: omit that label's summand and retain the divisor $m$, assuming
endpoint separation for all the remaining labels.  The constant $C$ may
depend on $b,\ell,C_{\rm reg}$, but is independent of small $t$.
\end{proposition}

\begin{proof}
We estimate the first leg from an outer endpoint $y_i$ to the random
midpoint $u$; reversing time gives the identical estimate for the leg
from $z_i$ to $u$.  On $\mathcal C_v$, the endpoint windows of both
paths have diameter at most $s_0/8$.  Endpoint separation consequently
leaves distance at least $3s_0/4>r$ throughout $[0,h_0]$.
Every mesh interval on which a collision remains possible starts at
$s\ge h_0-\tau\ge h_0/2$.

We first explain the use of lifts.  On a mesh interval, the obstacle has
diameter at most $rD^\gamma<1$.  A collision with a confined replacement
therefore requires $d(c(s),v)<12$ at the interval's starting time.  For
$K>50$ this determines a unique lift $\widetilde c(s)$ within distance
$12$ of the chosen lift of $v$; continue that lift over the interval.
Its displacement is at most $rA_k(c)$.  Intervals for which no such lift
exists make no contribution.  We may thus use Euclidean Gaussian
densities without summing over windings.

\paragraph{The first half of a leg.}
Suppose $s\le t/2$ and condition on $u$.  If $X_s$ is the position of
the trial bridge at time $s$, its segment of length $\tau$ has the
representation
\begin{equation}\label{bridge:markov-segment}
 X_{s+q}=X_s+\frac{q}{t-s}(u-X_s)+R_q,
 \qquad 0\le q\le\tau.
\end{equation}
Here $R$ is the initial segment of a centered bridge of duration $t-s$,
independent of $X_s$.  Its law does not depend on $u$ or on the endpoint
label.  This is the Gaussian conditional law of a bridge after its
time-$s$ position has been specified.  Equivalently,
$R_q=W_q-qW_{t-s}/(t-s)$ for a Brownian motion $W$ of generator $\Delta$.
Gaussian maximal estimates, or the Brownian reflection bound integrated
against $3a^2\dd a$, give
\begin{equation}\label{bridge:noise-moment}
 B_s=\frac1r\sup_{0\le q\le\tau}|R_q|,
 \qquad \E B_s^3\le C,
\end{equation}
because $\tau\le r^2$ and $\tau\le t-s$.

On the cutoff event, $|X_s-v|\le10$, whereas every target point is within
distance $2$ of $v$.  The drift term in
\eqref{bridge:markov-segment} is thus at most $C\tau/t\le r$.
A collision on this interval, together with the cutoff, implies
\begin{equation}\label{bridge:necessary-ball}
 |X_s-\widetilde c(s)|
 \le r\bigl(2+A_k(c)+B_s\bigr).
\end{equation}
We use the cutoff only to obtain this necessary event.  We now discard
it and integrate the unrestricted bridge law.  By
Lemma~\ref{bridge:gaussian-packing}, the density of $X_s$, averaged over
the labels, is bounded by \eqref{bridge:packed-density}.  Since $B_s$
is independent of $X_s$ and has the common bound
\eqref{bridge:noise-moment}, the averaged probability of
\eqref{bridge:necessary-ball} is at most
\begin{equation}\label{bridge:early-step}
 Cr^3(1+A_k(c))^3
 \left(1+\frac{1}{Ds^{3/2}}\right).
\end{equation}
This estimate holds for each $u$, hence also after averaging $u$ over
the target cell.  Omitting one label only decreases the unnormalized
Gaussian sum and leaves the same bound with divisor $m$.

\paragraph{The second half of a leg.}
Suppose $s>t/2$.  For a fixed label, write the whole unrestricted leg as
\[
 X_q=(1-q/t)y_i+(q/t)U+\beta_q,\qquad 0\le q\le t,
\]
where $U$ is uniform in $B$ and the centered bridge $\beta$ is
independent of $U$.  Condition on the entire path $\beta$.  The
conditional density of $X_s$ is at most $(t/s)^3\le8$.
The mean displacement during the mesh interval is at most
$C\tau/t\le r$, uniformly in $U\in B$.  Put
\[
 \widetilde B_s=
 r^{-1}\sup_{0\le q\le\tau}|\beta_{s+q}-\beta_s|.
\]
Writing $\beta_q=W_q-qW_t/t$ gives
$\E\widetilde B_s^3\le C$: the Brownian increment contributes
$C\tau^{3/2}/r^3$, and the linear correction contributes
$C\tau^3/(r^3t^{3/2})$, both bounded.  Conditional on $\beta$, a
collision requires the ball in \eqref{bridge:necessary-ball}, with
$B_s$ replaced by $\widetilde B_s$.  Both its radius and center are then
fixed while $U$ retains its uniform distribution.  Integrating first
over $U$ and then over $\beta$ bounds the probability by
\begin{equation}\label{bridge:late-step}
 Cr^3(1+A_k(c))^3.
\end{equation}
In particular, no independence between $X_s$ and the fluctuation
increment is needed here.

\paragraph{Summing the intervals.}
Apply the two bounds to both legs.  The terms without $s^{-3/2}/D$
use the averaged regularity condition:
\[
 Cr^3\sum_{k=1}^{2m_t}(1+A_k(c))^3
 \le Cr^3m_t\le Ctr.
\]
Here $m_t=t/\tau$ and $\tau\ge r^2/2$.  For the remaining terms use
$(1+A_k(c))^3\le CD^{3\gamma}$ and
\[
 \sum_{\substack{s=j\tau\ge h_0/2\\s\le t/2}}s^{-3/2}
 \le \frac{C}{\tau\sqrt{h_0}}
 \le \frac{C}{r^2\sqrt{h_0}}.
\]
Their contribution is therefore at most
$Cr^3D^{3\gamma}/(Dr^2\sqrt{h_0})
=CrD^{3\gamma}/(D\sqrt{h_0})$.
These two sums give \eqref{bridge:averaged-estimate}.  Every constant in these
sums is independent of $t$; making $t$ smaller only increases the
required lower threshold on $D$ in
\eqref{bridge:scale-restrictions}.
\end{proof}

\begin{proposition}[A single replacement]\label{bridge:single-hit}
Let $y\in B_v$, $|z-v|\le3$, and let $c$ be a deterministic regular
path satisfying $d(y,c(0))>s_0$ and $d(z,c(2t))>s_0$.
For either $Q=Q_{y,z}^{B}$, with any target cell $B$, or
$Q=Q_{y,z}^{0}$, one has
\begin{equation}\label{bridge:single-estimate}
 Q(\mathcal C_v\cap\mathcal H(c))
 \le Ctr+C\frac{rD^{3\gamma}}{\sqrt{h_0}}
 \le C\alpha\bigl(tD^{-1}+D^{-1/4}\bigr).
\end{equation}
Thus this probability tends to zero with $D$ for each fixed $t$.
\end{proposition}

\begin{proof}
For $Q_{y,z}^{B}$, repeat the preceding proof with one endpoint.  At
$s\le t/2$ its Gaussian density is at most $Cs^{-3/2}$; the bound
\eqref{bridge:early-step} becomes
$Cr^3(1+A_k(c))^3s^{-3/2}$.  The second half of each leg retains
\eqref{bridge:late-step}.  The same two sums give the first inequality.

For $Q_{y,z}^{0}$, use the forward orientation up to time $t$ and the
reversed orientation for the rest.  In either orientation the
remaining bridge duration is at least $t$.  Formula
\eqref{bridge:markov-segment} holds with $2t$ in place of $t$, its
drift on the cutoff event is at most $C\tau/t$, and its time-$s$
density is bounded by $Cs^{-3/2}$ for $0<s\le t$.  Endpoint separation
again excludes the window $[0,h_0]$.  Consequently the first-half
argument and its mesh sum apply over both orientations, giving
$CrD^{3\gamma}/\sqrt{h_0}$.  Adding the nonnegative term $Ctr$ yields
the stated common bound.  Finally $r=\alpha/D$, $h_0=D^{-1}$ and
$3\gamma=1/4$ give the last inequality.
\end{proof}

\begin{remark}\label{bridge:midpoint-quantifier}
The midpoint in $Q_{y,z}^{B}$ is averaged over a unit cell.  An estimate
uniform in a prescribed midpoint is false: if $u=c(t)$, the trial
collides at time $t$.  Restricting a uniform midpoint to a measurable
set of volume at least $a>0$ multiplies an upper bound for a nonnegative
event by at most $a^{-1}$.  Conditioning a trial on an event of
probability at least $a$ has the same cost.  These elementary changes
of measure will allow us to use
Proposition~\ref{bridge:single-hit} after choosing successful trials.
\end{remark}

\subsection{Many paths can be moved successfully}

The averaged estimate now pays for all nearby obstacles.  The following
statement separates that deterministic estimate from the random cell
certificates that will supply its hypotheses.

\begin{corollary}[Supply of successful labels]\label{bridge:success-supply}
Consider a collection of deterministic torus paths with endpoints
$(y_j,z_j)$, and fix a cell $v$.  Suppose that at most $H_1D$ of these
paths visit the ball of radius $12$ about $v$, and that every visiting
path is regular.  Suppose a list of $m\ge\ell D/2$ labels satisfies
$y_i\in B_v$, $|z_i-v|\le3$, and that both endpoints of every listed
label are at distance greater than $s_0$ from every other endpoint in
the corresponding full configuration.

Call a trial for label $i$ successful if it satisfies $\mathcal C_v$
and avoids every obstacle path except the original path of $i$.
For sufficiently small fixed $t>0$ and then sufficiently large $D$,
at least $9m/10$ listed labels have the following property: for each
of the seven target cells $B$, the set
\begin{equation}\label{bridge:successful-target-set}
 T_{i,B}=\{u\in B:
     Q_{y_i,z_i}^u(\text{successful trial})\ge9/10\}
\end{equation}
has volume at least $99/100$.
All choices are uniform in the deterministic collection, $v$, and
$\alpha\in[\alpha_0/2,2\alpha_0]$.
\end{corollary}

\begin{proof}
The linear means of both bridge legs stay within distance $3$ of $v$.
The Gaussian maximal bound therefore makes confinement failure at
most $Ce^{-c/t}$.  Lemma~\ref{path:regularity} gives a regularity
failure probability $\varepsilon_D(t)\to0$, uniformly over the listed
endpoints and the target points.  A confined trial can collide only
with an obstacle visiting the ball of radius $12$, since $r<1$.
Union-bounding over those obstacles and using
Proposition~\ref{bridge:averaged-hit}, with a self-label summand
omitted when necessary, gives for each target cell
\begin{align}
 \frac1m\sum_i\int_B
 Q_{y_i,z_i}^{u}(\text{failure})\dd u
 &\le Ce^{-c/t}+\varepsilon_D(t)
       +CH_1D\left(tr+\frac{rD^{3\gamma}}{D\sqrt{h_0}}\right)
 \notag\\
 &\le Ce^{-c/t}+\varepsilon_D(t)
       +CH_1\alpha\bigl(t+D^{-1/4}\bigr).
 \label{bridge:total-failure}
\end{align}
The constant multiplying $t$ is independent of small $t$, which is
the reason for the two parts of the bridge estimate.  First choose
$t$ small and then $D$ large so that the last bound is less than
$1/70000$ for every $\alpha$ in the fixed window.

If $|T_{i,B}|<99/100$, the integral of the failure probability over
$B$ is greater than $(1/100)(1/10)=1/1000$.  Thus fewer than
$m/70$ labels fail the conclusion for a specified target cell.
There are seven target cells, so fewer than $m/10$ labels fail for
at least one of them.  All the remaining labels have the asserted
property.
\end{proof}
